% === IMÁGENES: limitar al ancho de página sin agrandar ===
\usepackage[export]{adjustbox}
\let\oldincludegraphics\includegraphics
\renewcommand{\includegraphics}[2][]{\oldincludegraphics[max width=\linewidth,max height=0.4\textheight,keepaspectratio,#1]{#2}}

% === CÓDIGO: word-wrap, números de línea DENTRO del recuadro ===
\usepackage{fvextra}
\usepackage{xcolor}
\definecolor{codebg}{HTML}{f5f5f5}
\definecolor{codeframe}{HTML}{e0e0e0}
\definecolor{linenumcolor}{HTML}{999999}

\DefineVerbatimEnvironment{Highlighting}{Verbatim}{
  breaklines=true,
  breakanywhere=true,
  fontsize=\footnotesize,
  numbers=left,
  numbersep=10pt,
  xleftmargin=15pt,
  commandchars=\\\{\}
}

% Redefinir Shaded para un recuadro que incluya los números
\usepackage{framed}
\renewenvironment{Shaded}{%
  \definecolor{shadecolor}{HTML}{f5f5f5}%
  \setlength{\OuterFrameSep}{0pt}%
  \begin{snugshade*}%
}{%
  \end{snugshade*}%
}

% Color de números de línea
\renewcommand{\theFancyVerbLine}{\textcolor{linenumcolor}{\scriptsize\arabic{FancyVerbLine}}}

% === TABLAS: ajustar al ancho de página ===
\usepackage{longtable}
\usepackage{booktabs}
\usepackage{array}
\usepackage{calc}
\let\oldlongtable\longtable
\renewcommand{\longtable}{\footnotesize\oldlongtable}
\renewcommand{\arraystretch}{1.4}

% === Evitar cortes en bloques de código ===
\usepackage{float}
\floatplacement{figure}{H}

% === BLOCKQUOTES y ADMONITIONS con colores por tipo ===
\usepackage{mdframed}
\usepackage{xcolor}

\definecolor{adm-note}{HTML}{1a73e8}
\definecolor{adm-tip}{HTML}{0d904f}
\definecolor{adm-warning}{HTML}{e65100}
\definecolor{adm-danger}{HTML}{d32f2f}
\definecolor{adm-info}{HTML}{0288d1}
\definecolor{adm-example}{HTML}{7b1fa2}

% Estilo por defecto (se redefine antes de cada admonition)
\global\mdfdefinestyle{admonitionstyle}{
  linewidth=2pt,
  linecolor=adm-note,
  backgroundcolor=adm-note!5,
  leftmargin=0pt, rightmargin=0pt,
  innerleftmargin=10pt, innerrightmargin=10pt,
  innertopmargin=8pt, innerbottommargin=8pt,
  skipabove=10pt, skipbelow=10pt,
  topline=false, rightline=false, bottomline=false,
  nobreak=true
}
\renewenvironment{quote}{\begin{mdframed}[style=admonitionstyle]}{\end{mdframed}}

% === TÍTULOS: numerados y con más peso visual ===
\usepackage{titlesec}
\setcounter{secnumdepth}{5}
\setcounter{tocdepth}{4}

\titleformat{\chapter}[display]
  {\normalfont\Huge\bfseries}
  {\chaptertitlename\ \thechapter}{20pt}{\Huge}
\titlespacing*{\chapter}{0pt}{-30pt}{30pt}

\titleformat{\section}
  {\normalfont\LARGE\bfseries}{\thesection}{1em}{}
\titlespacing*{\section}{0pt}{20pt}{10pt}

\titleformat{\subsection}
  {\normalfont\Large\bfseries}{\thesubsection}{1em}{}
\titlespacing*{\subsection}{0pt}{15pt}{8pt}

\titleformat{\subsubsection}
  {\normalfont\large\bfseries}{\thesubsubsection}{1em}{}
\titlespacing*{\subsubsection}{0pt}{12pt}{6pt}

\titleformat{\paragraph}[hang]
  {\normalfont\normalsize\bfseries}{\theparagraph}{1em}{}
\titlespacing*{\paragraph}{0pt}{8pt}{4pt}

\titleformat{\subparagraph}[hang]
  {\normalfont\normalsize\bfseries}{\thesubparagraph}{1em}{}
\titlespacing*{\subparagraph}{0pt}{6pt}{3pt}

% === PÁRRAFOS ===
\usepackage{parskip}
\setlength{\parskip}{0.4em}
\setlength{\parindent}{0pt}

% === LINKS ===
\usepackage{hyperref}
\hypersetup{colorlinks=true, linkcolor=blue!60!black, urlcolor=blue!60!black}

% === CENTRAR figcaptions ===
\usepackage{caption}
\captionsetup{justification=centering}
